\scanpage{86}
donné. On a dans ce cas
\[ \left\{\begin{aligned}u(x,y)&=(1-y,-x)\\
u(x,y)&=(x+1,y-1)\end{aligned}\right. \tag*{(25)} \]
i.e. $u^2$ est la translation par le vecteur $(1,-1)$.
Donc en car. $0$ $u^2$ est d'ordre infini, donc
a fortiori $u$ l'est, donc la réalisation obtenue
de $\Pi_\infty$ est d'ordre infini (on trouve
qu'elle est fidèle\ldots) En car. $p>0$, $u^2$ est
d'ordre $p$, donc $u$ d'ordre $2p$, on trouve
\deleted{\unsure{un $2p$|un $2p$-}} donc que la réalisation géom impropre
se factorise par $\Pi_{2p}$. On vérifie
aisément que la réalisation géom de $\Pi_{2p}$
est bel et bien fidèle, donc donne un
$2p$-polygône régulier.

Avec la théorie projective des polygônes
réguliers, l'anomalie \unsure{apparente} des ``polygônes
réguliers impropres'' s'explique complètement, et on voit que du point de vue
projectif on n'a rien perdu en excluant
le cas impropre, ce qu'on fera désormais.

\uline{Scholie} Le cas d'une réalisation géom
\margin{terminologie ``invariant numérique''}
régulière \uline{propre} de $\Pi_\infty$ est classé par $\alpha\in k$,
qu'on peut choisir arbitrairement\ldots
% archive sheet number: 85
